\documentclass[11pt,a4paper]{article}

% ---------- Codifica e lingua ----------
\usepackage[T1]{fontenc}
\usepackage[utf8]{inputenc}
\usepackage[main=italian,provide=*]{babel}

% ---------- Matematica ----------
\usepackage{amsmath,amssymb}
\usepackage{mathtools}
\usepackage{xcolor}
\definecolor{amber}{HTML}{C68A00}
\definecolor{teal2}{HTML}{2E8B7A}
\usepackage{tikz}

% ---------- Impaginazione ----------
\usepackage[a4paper,margin=2.7cm]{geometry}
\usepackage{caption}
\usepackage{booktabs}
\usepackage{enumitem}
\usepackage{placeins}
\usepackage[colorlinks=true,linkcolor=black,citecolor=black,urlcolor=blue!55!black]{hyperref}

\newcommand{\ket}[1]{\lvert #1 \rangle}
\newcommand{\bra}[1]{\langle #1 \rvert}
\newcommand{\braket}[1]{\langle #1 \rangle}
\newcommand{\Tr}{\operatorname{Tr}}

\title{NOTA DI LAVORO INTERNA --- non per la commissione\\
\large Derivazione e verifica dell'argomento di continuità\\
\large(VQE noise-aware: perché $N^*$ non si sposta)}
\author{Samuele \\ Relatore: Prof.\ Alessandro Chiesa}
\date{}

\begin{document}
\maketitle

\begin{center}
\fbox{\parbox{0.9\linewidth}{\centering\textbf{Questo documento è una
verifica metodologica interna, non un capitolo della tesi.} Il risultato
--- l'argomento di continuità è confermato, con un residuo del primo
ordine inferiore all'$1\%$ --- va citato in tesi in una riga; questa nota
esiste per essere sicuri di quella riga, non per essere trasformata in
sezione.}}
\end{center}

\tableofcontents

\section{Il problema}

Nei documenti dei risultati (\texttt{risultati\_vqe\_noise\_aware\_definitivo.tex},
§\,``Perché il risultato generalizza: l'argomento di continuità'') si afferma:

\begin{quote}
$\theta^*_\text{noise-aware}$ è quasi indistinguibile da
$\theta^*_\text{ideale}$; qualunque funzione ragionevolmente liscia dello
stato preparato deve differire fra le due preparazioni di una quantità
altrettanto minuscola, per continuità.
\end{quote}

Questa è un'affermazione qualitativa, non dimostrata. Questo documento la
rende quantitativa: deriva un legame esplicito fra $\|\Delta\theta\|$ e lo
scostamento osservato nelle due metriche a valle (fedeltà di Trotter,
modulo del correlatore), e lo confronta col dato numerico già raccolto.

\section{Derivazione}

\subsection{Linearizzazione al primo ordine}

Sia $f:\mathbb R^3\to\mathbb R$ una qualunque delle due metriche a valle
(fedeltà $F(N;\theta)$ a $N$ fissato, o modulo del correlatore
$|C(N;\theta)|$), vista come funzione dei tre parametri variazionali
$\theta=(\theta_0,\theta_1,\theta_2)$ dell'ansatz PMA-2q$\cdot$3, a $N$ e
a tutti gli altri parametri fisici fissati. Se $f$ è differenziabile in
un intorno di $\theta^*_\text{ideale}$, lo sviluppo di Taylor al prim'ordine dà
\begin{equation}
f(\theta^*_\text{ideale}+\Delta\theta) = f(\theta^*_\text{ideale}) +
\nabla f(\theta^*_\text{ideale})\cdot\Delta\theta + O(\|\Delta\theta\|^2),
\label{eq:taylor}
\end{equation}
dove $\Delta\theta\equiv\theta^*_\text{noise-aware}-\theta^*_\text{ideale}$.
Per la disuguaglianza di Cauchy--Schwarz,
\begin{equation}
|\Delta f| \le \|\nabla f(\theta^*_\text{ideale})\|\;\|\Delta\theta\| +
O(\|\Delta\theta\|^2).
\label{eq:bound}
\end{equation}
Questo è il legame quantitativo che l'affermazione originale lasciava
implicito: lo scostamento sull'osservabile è limitato dal prodotto fra la
norma del gradiente e la norma dello scostamento dei parametri --- non
``piccolo perché sì'', ma piccolo con un fattore di proporzionalità
preciso e calcolabile.

\subsection{$f$ è davvero liscia? Una cautela necessaria}
\label{sec:cautela-liscezza}

L'Eq.~\eqref{eq:taylor} presuppone che $f(\theta)$ sia differenziabile.
Fisicamente lo è: ogni stato $\ket{\psi(\theta)}$ preparato dall'ansatz è
combinazione di seni e coseni di $\theta/2$ (rotazioni $R_y$), il canale
di rumore agisce linearmente sull'operatore densità, e la fedeltà/il
modulo del correlatore sono funzioni polinomiali degli elementi di
matrice di $\rho(\theta)$ --- quindi $f(\theta)$ è $C^\infty$, anzi
analitica reale, in $\theta$.

\textbf{Ma la funzione calcolata dal codice non è necessariamente quella
funzione.} Il circuito usato per calcolare $f(\theta)$ nella pipeline
principale (\texttt{trotter\_rumoroso\_dimero.py},
\texttt{correlatori\_rumorosi\_dimero.py}) transpila l'ansatz \emph{dopo}
aver assegnato il valore numerico di $\theta$ --- la scelta corretta per
il conteggio di gate minimo (già stabilita in una sessione precedente),
ma con una conseguenza collaterale qui rilevante: il transpilatore può
scegliere sintesi diverse dei gate a seconda del valore numerico
specifico di $\theta$ (per esempio se un angolo è vicino a $0$ o $\pi$,
alcune rotazioni si fondono o si cancellano). Questo rende la funzione
\emph{numericamente calcolata} discontinua a scala di transpilazione,
anche se la funzione \emph{fisica} sottostante è liscia.

\begin{center}
\renewcommand{\arraystretch}{1.3}
\begin{tabular}{@{}ccc@{}}
\toprule
$h$ (passo differenze finite) & ops($\theta_2+h$) & ops($\theta_2-h$)\\
\midrule
$10^{-3}$ & 12 \texttt{rz} & 11 \texttt{rz}\\
$10^{-5}$ & 12 \texttt{rz} & 11 \texttt{rz}\\
$10^{-7}$ & 11 \texttt{rz} & 11 \texttt{rz}\\
\bottomrule
\end{tabular}
\end{center}

\noindent Verificato direttamente (tabella sopra): il conteggio dei gate
\texttt{rz} del circuito transpilato cambia fra le due valutazioni usate
per stimare la derivata rispetto a $\theta_2$, per $h=10^{-3}$ e
$h=10^{-5}$. Conseguenza: una stima del gradiente per differenze finite
calcolata direttamente sulla pipeline principale è \textbf{instabile in
modo drammatico} al variare di $h$:

\begin{center}
\renewcommand{\arraystretch}{1.2}
\begin{tabular}{@{}cccc@{}}
\toprule
$h$ & $\partial F/\partial\theta_0$ & $\partial F/\partial\theta_1$ & $\partial F/\partial\theta_2$\\
\midrule
$10^{-3}$ & $0.2309$ & $0.0594$ & $-0.169$\\
$10^{-4}$ & $1.329$ & $-1.039$ & $-0.059$\\
$10^{-5}$ & $0.2309$ & $0.0594$ & $\mathbf{-11.04}$\\
$10^{-6}$ & $0.2309$ & $0.0594$ & $-0.0594$\\
$10^{-7}$ & $\mathbf{1099}$ & $0.0594$ & $-0.0594$\\
\bottomrule
\end{tabular}
\end{center}

\noindent Non è rumore fisico: è un artefatto della transpilazione, e
avrebbe reso qualunque stima del gradiente fatta ingenuamente sulla
pipeline principale inaffidabile (nella prima prova, con $h=10^{-5}$, la
componente $\theta_2$ del gradiente usciva 190 volte più grande delle
altre due --- un artefatto, non un fatto fisico).

\subsection{La correzione: stessa idea della trappola di transpilazione, usata al contrario}

La soluzione è la stessa tecnica già nota (transpilare l'ansatz
\textbf{una sola volta}, in forma parametrica/simbolica, e poi assegnare
i valori numerici) --- ma qui usata per lo scopo \emph{opposto} a quello
per cui era stata scartata nell'ottimizzazione VQE noise-aware:
\begin{itemize}[leftmargin=1.3em,itemsep=0.2em]
\item nel VQE con termine DM sotto rumore, la transpilazione simbolica dà un
circuito \emph{peggiore} (più rumoroso, più gate) --- per quello si
transpila dopo aver assegnato i parametri, ad ogni valutazione;
\item qui, per stimare una \textbf{derivata}, serve che la struttura del
circuito resti \emph{esattamente la stessa} al variare di $\theta$ in un
intorno infinitesimo --- per quello la transpilazione simbolica (fatta
una volta sola, poi solo \texttt{assign\_parameters}) è l'unica scelta
corretta, anche se il circuito risultante ha più gate (19 \texttt{rz} + 14
\texttt{sx} invece di 11+8) e quindi una fedeltà assoluta leggermente più
bassa a parità di rumore.
\end{itemize}
Con questa correzione, la stima del gradiente per differenze finite è
\textbf{stabile su quattro ordini di grandezza di $h$} (verificato,
$h=10^{-3}\ldots10^{-6}$, differenza $<10^{-6}$ sulle componenti del
gradiente) --- la firma di una derivata numerica corretta.

\section{Verifica numerica}

\subsection{Impostazione}

Usando la costruzione a struttura fissa (Sez.~\ref{sec:cautela-liscezza}),
al punto di lavoro ``test 2'' ($b/J=0.35$, $D/J=0.80$), con
$\theta^*_\text{ideale}$ e $\theta^*_\text{noise-aware}$ registrati nelle
sessioni precedenti:
\begin{equation}
\Delta\theta = \theta^*_\text{noise-aware}-\theta^*_\text{ideale} =
(1.993, -1.278, -0.559)\times10^{-4}, \qquad \|\Delta\theta\| =
2.433\times10^{-4}.
\end{equation}

\subsection{Fedeltà di Trotter, $N=8$}

Gradiente calcolato per differenze finite (struttura fissa, $h=10^{-5}$,
stabile su 4 ordini di grandezza di $h$):
\begin{equation}
\nabla F = (0.22945,\ 0.05892,\ -0.05892), \qquad \|\nabla F\| = 0.2441.
\end{equation}

\begin{center}
\renewcommand{\arraystretch}{1.3}
\begin{tabular}{@{}lc@{}}
\toprule
\textbf{Quantità} & \textbf{Valore}\\
\midrule
$F(\theta^*_\text{ideale})$ & $0.813472$\\
$F(\theta^*_\text{noise-aware})$ & $0.813514$\\
$\Delta F$ osservato & $4.147\times10^{-5}$\\
$\Delta F$ predetto ($\nabla F\cdot\Delta\theta$) & $4.150\times10^{-5}$\\
rapporto predetto/osservato & $1.0007$\\
residuo (ordine $\ge2$) & $-3.1\times10^{-8}$\\
\bottomrule
\end{tabular}
\end{center}

\subsection{Modulo del correlatore, $N=5$}

\begin{equation}
\nabla|C| = (-0.3167,\ 0.1104,\ -0.1663), \qquad \|\nabla|C|\| = 0.3744.
\end{equation}

\begin{center}
\renewcommand{\arraystretch}{1.3}
\begin{tabular}{@{}lc@{}}
\toprule
\textbf{Quantità} & \textbf{Valore}\\
\midrule
$|C(\theta^*_\text{ideale})|$ & $0.248044$\\
$|C(\theta^*_\text{noise-aware})|$ & $0.247976$\\
$\Delta|C|$ osservato & $-6.792\times10^{-5}$\\
$\Delta|C|$ predetto ($\nabla|C|\cdot\Delta\theta$) & $-6.794\times10^{-5}$\\
rapporto predetto/osservato & $1.0004$\\
residuo (ordine $\ge2$) & $+2.4\times10^{-8}$\\
\bottomrule
\end{tabular}
\end{center}

\begin{center}
\fbox{\parbox{0.85\linewidth}{\centering
In entrambi i casi la previsione lineare riproduce lo scostamento
osservato entro lo $0.1\%$: siamo saldamente nel regime dove il
prim'ordine domina, il residuo di ordine superiore è tre ordini di
grandezza più piccolo del termine lineare stesso.}}
\end{center}

\section{Cosa dice, e cosa non dice, questo risultato}

\begin{itemize}[leftmargin=1.3em,itemsep=0.3em]
\item \textbf{Confermato}: al punto ``test 2'', per entrambe le metriche
controllate, lo scostamento osservato è spiegato quasi esattamente
(entro l'$1\%$) dal termine lineare $\nabla f\cdot\Delta\theta$. Non è
più un'affermazione qualitativa (``dovrebbe essere piccolo'') ma un
conto verificato (``è esattamente piccolo quanto la teoria del
prim'ordine predice'').
\item \textbf{Non dimostrato in generale}: il bound
dell'Eq.~\eqref{eq:bound} usa $\|\nabla f(\theta^*_\text{ideale})\|$
calcolato in questo punto specifico. Se in un altro punto di lavoro (o
per un'altra metrica) il gradiente fosse molto più grande, lo stesso
$\|\Delta\theta\|$ produrrebbe uno scostamento proporzionalmente più
grande --- il bound protegge \emph{a patto che} il gradiente resti
moderato, cosa verificata qui ($\|\nabla F\|\sim0.24$,
$\|\nabla|C|\|\sim0.37$, entrambi ordine 1) ma non garantita a priori per
ogni punto.
\item \textbf{La cautela più importante}: qualunque stima futura di
derivate o sensibilità su questa pipeline deve usare la costruzione a
struttura fissa (transpilazione simbolica una tantum), mai la
transpilazione post-assegnazione usata per l'efficienza --- la Sez.
\ref{sec:cautela-liscezza} mostra cosa succede altrimenti (un artefatto
del tutto scollegato dalla fisica, capace di produrre un gradiente
apparente 190 volte più grande del vero).
\end{itemize}

\section{Riga da riportare in tesi, se serve}

Se si vuole citare questo risultato senza riportare l'intera derivazione:
\begin{quote}
\emph{L'argomento di continuità è stato verificato quantitativamente: lo
scostamento osservato in entrambe le metriche a valle è riprodotto entro
lo 0.1\% da una previsione lineare $\nabla f\cdot\Delta\theta$, confermando
che il regime è dominato dal prim'ordine (nota di lavoro interna, non
riportata per esteso).}
\end{quote}

\end{document}
